\chapter*{Abstract}
\addcontentsline{toc}{chapter}{Abstract}

HTTP/3 uses QUIC/UDP to reduce latency and improve resilience on poor networks. Node.js does not offer stable native support, so this thesis asks whether HTTP/3 can work in Node.js REST APIs, what obstacles arise, how it performs compared to HTTP/1.1 and HTTP/2, and how to make it work in practice. We combined published research, a literature review, and experiments with a real API. We built an API for load testing with three routes: HTTP/1.1 directly, and HTTP/2 and HTTP/3 through an NGINX reverse proxy. We used k6 to load test HTTP/1.1 and HTTP/2, and a QUIC client in Docker for HTTP/3. We ran two sets of tests: one with the network degraded (using tc netem, GET /bench/ping, 50 virtual users across four network profiles) and one with a clean network (different payload sizes, echo upload, and 10, 50, or 100 simultaneous users). Each test ran ten times, with results summarized by median and interquartile range. HTTP/3 ran without failures across all tests (120 runs with degraded network, 210 with clean network), which shows it works using a proxy at the edge while the backend stays on HTTP/1.1. However, different load generators, asymmetric topologies (HTTP/1.1 direct versus HTTP/2 and HTTP/3 via NGINX), and QUIC client limits made general conclusions difficult. These results describe this specific setup. The proxy approach required no code changes in the application, but Alt-Svc negotiation and protocol fallback were not tested. The code is open and the collection scripts at \url{https://github.com/mendeslian/QUICBench-Node} allow others to replicate and extend the tests.

\noindent\textbf{Key-words:} HTTP/3. QUIC. Node.js. REST APIs. Feasibility; Performance.
